\section{The Memoirs and the Textual History of the Secret}\label{sec:memoirs}

\subsection{Late composition, and why it matters}

Nearly everything the devout world knows as ``the message of
F\'atima'' beyond the bare 1917 record comes from documents Lucia
wrote between eighteen and twenty-seven years after the events, in
obedience to her bishop: the First Memoir (finished 25 December 1935),
centered on Jacinta; the Second (21 November 1937), which first
narrates the 1916 angel apparitions at length and her own interior
history; the Third (31 August 1941), which first reveals the first
two parts of the secret; and the Fourth (8 December 1941), a
comprehensive account with annotations---including her note that,
except the part of the secret she was not then permitted to reveal,
``I shall say everything. I shall not knowingly omit anything, though
I suppose I may forget just a few small details of minor
importance.''\endnote{Memoir dates: the Sanctuary's seer chronology
(25 December 1935; 21 November 1937; 31 August 1941; 8 December
1941), corroborated for the Third and Fourth by the CDF's
\work{Message of Fatima} (2000), which prints the secret's first two
parts from the Third Memoir and quotes the Fourth Memoir's
undertaking in its note 6. The Dicastery for the Causes of Saints'
2023 decree likewise dates the memoirs 1935--1941, ``by order of the
Bishop of Leiria.'' Fifth and Sixth Memoirs (1989, 1993) concern her
parents and are not used here.}

The historiographical care this demands is not suspicion but
proportion. The memoirs are retrospective spiritual autobiography,
written after the Russian revolutions of 1917, after the rise of
Soviet communism and its persecutions, after the great aurora of
25 January 1938, and---for the Third and Fourth---into the Second
World War; their author herself distinguishes what she may reveal
from what she may not, and warns of forgetfulness in details. The
1917 documents verify that a secret existed and that a miracle was
promised three months ahead; they do not contain the 1941 wording.
Sound method therefore reads each memoir as a witness to what Lucia,
in 1935--1941, believed and remembered herself to have experienced
---a witness the Church took seriously enough to build on, but not a
stenographic deposit of 1917. The Church's own 2000 commentary
adopted precisely this frame for the visions themselves
(\S\ref{sec:interp2000}, \S\ref{sec:mode}).

\subsection{The three parts: composition, custody, publication}

\textbf{Parts one and two.} These are the first part, the vision of
hell shown to the children on 13 July 1917, and the second, the
announcement of devotion to the Immaculate Heart of Mary with its
conditional prophecy concerning Russia and the two wars, both as
written in the Third Memoir of 31 August 1941 and annotated in the
Fourth of 8 December 1941. Their wording is printed entire, with the
two Fourth Memoir annotations, in \S\ref{sec:text}; here only their
publication history is at issue. These two parts entered public
circulation in the early 1940s---the standard dating is 1942, in
anniversary publications authorized in Portugal---so that by the
1950s they were common property of theological discussion; the 2000
edition simply describes them as ``already\ldots\ published and
therefore known.'' The 1942 first imprints were not re-verified here
(Appendix~A).\endnote{CDF,
\work{Message of Fatima}, Bertone introduction; the theological
commentary cites the state of discussion in E.~Dhanis, ``Sguardo su
Fatima e bilancio di una discussione,'' \work{La Civilt\`a Cattolica}
104 (1953) II, 392--406. Pius XII's F\'atima-jubilee radio message of
31 October 1942 (\S\ref{sec:consecration}) belongs to the same anniversary moment.}

\textbf{Part three.} Lucia wrote the third part at Tuy on 3 January
1944, ``by order of His Excellency the Bishop of Leiria and the Most
Holy Mother\ldots,'' after months of anguished hesitation; there is
one manuscript, which the 2000 edition reproduces photostatically.
The sealed envelope stayed with the Bishop of Leiria until, ``to
ensure better protection,'' it was placed in the Secret Archives of
the Holy Office on 4 April 1957. John XXIII received the envelope on
17 August 1959, said---as the archives record---``We shall wait. I
shall pray. I shall let you know what I decide,'' and decided to
return it unpublished; Paul VI read the text on 27 March 1965 and
likewise returned it; John Paul II asked for it after the attempt on
his life of 13 May 1981.\endnote{All custody facts: CDF,
\work{Message of Fatima}, Bertone introduction, with its notes 1--2
(John XXIII's diary for 17 August 1959; the July--August 1981
movements of the two envelopes). The envelope's text and the
translation policy---``the original text\ldots\ has been respected,
even as regards the imprecise punctuation''---are the edition's.}

\textbf{The 1960 condition.} On the outer envelope Lucia had written
that it could be opened only after 1960, and only by the Patriarch of
Lisbon or the Bishop of Leiria. Decades of speculation attached to
that year; her own recorded explanation of it, and what the choice
does and does not imply, are taken up in \S\ref{sec:1960}.

\textbf{Publication, 26 June 2000.} At the end of the beatification
Mass at F\'atima on 13 May 2000, Cardinal Angelo Sodano, Secretary of
State, announced by the Pope's direction that the third part would be
published after preparation of a commentary. The Congregation for the
Doctrine of the Faith published \work{The Message of Fatima} on
26 June 2000: a documentary dossier containing Bertone's
introduction; the photostat and translation of the first two parts
(Third Memoir text with Fourth Memoir annotations); the photostat and
translation of the third part; John Paul II's letter of 19 April 2000
to Sister Lucia; the account of the 27 April 2000 conversation;
Sodano's 13 May announcement; and Cardinal Joseph Ratzinger's
theological commentary. The texts themselves are printed in
\S\ref{sec:text}; the Holy See's interpretation of them is
\S\ref{sec:interp2000}.\endnote{CDF, \work{Message of Fatima},
contents as listed on the Holy See's page; Sodano announcement of
13 May 2000, ibid. Lucia confirmed the manuscript and her handwriting
in the 27 April 2000 conversation.}

The completeness of that publication was later challenged in print;
the challenge, the Holy See's 2016 response, and this study's
assessment of the evidence are \S\ref{sec:withheld}. Textually, the
position of record is the 2000 edition's: one manuscript, reproduced
in facsimile, published entire.
